\documentclass{article}
\usepackage{amsmath,url}
\title{A retained-receipt example where two risk doctrines choose different probes}
\author{}\date{}
\begin{document}\maketitle
\paragraph{Stakeholder law and native cut.}
The old delivery is mandatory and completes in each replayed
branch. A next optional admission is rewarded only for an actual
successful callback; an attempted admission without that callback
is the adverse event. Paying for an informative probe uses genuine
receiver credit and also has a stipulated objective price $k\ge0$
per probed trial. No market or measured energy objective is implied.
There is one physical choice to run the authorized preceding probe
or omit it; the choice cannot depend on the hidden binding or RF
drop. Native no-probe and charged-probe histories are replayed
through the same selected fixture as the earlier receipt study.

Let $s$ be the fixed mass of a shared receiver; $d$ the fixed drop
marginal for the probe if tried. Both are stipulated, not estimated
from the native rows. Define $r$ as the mass of a separate receiver
with no drop, hence of a positive receipt. Complete laws allow exactly
\[
             L=\max(0,1-s-d)\le r\le U=\min(1-s,1-d).
\]
The four atom masses follow by marginal accounting and nonnegativity.
Each law is fixed for the entire decision, not resampled after a report.
For a two-byte next duty, all separate histories succeed, including
the separately charged positive case; shared histories fail.
For a three-byte next duty, the separate positive case instead
fails because its probe used the last receiver credit. Omitting
the probe makes either size feasible on the separate receiver.
All of these are native callbacks under the explicit authored path law.

\paragraph{Whole-plan cap (two-byte next duty).}
Let $l$ be the independent probe fraction, $n$ the probability
of admission jointly with omitting the probe, $a$ the admission
fraction weighted by choosing probe on positive response, and $b$
the corresponding weighted fraction on absence. Thus
$n\le1-l$, $a,b\le l$, all nonnegative; $a,b$ do not already include
the probability of the respective signal.
The missed-admission risk is $s(n+b)$, and good service is
\[
        (1-s)n+r a+(1-s-r)b .
\]
For cap $e$ let $x=\min(1,e/s)$ when $s>0$. Set $a=l$,
use $n=\min(1-l,x)$ and then $b=x-n$. This policy attains
$(1-s)x+L\min(l,1-x)$ in the worst law.
For any policy evaluate at $r=L$. Splitting the objective as
$(1-s)(n+b)+r(a-b)$ gives the bound $(1-s)x+r l$.
Splitting it as $(1-s-r)(n+b)+r(n+a)$ gives the bound
$(1-s)x+r(1-x)$, since $n+a\le1$. Hence the envelope is exact.
Subtracting $k l$ gives $l=0$ if $L\le k$, otherwise the smallest
optimal fraction is $1-x$. If the binding is always separate
($s=0$), no probe is necessary. If always shared, optional
service gives no benefit and can be withheld.

\paragraph{Risk at each visible decision history.}
Here the cap is imposed conditional on the probe choice and,
where relevant, the actual receipt, but before the independent
admission coin. It is not a cap conditional on an admission
already being made. The no-probe term obeys $n\le (1-l)x$.
On absence, $sb\le e l (1-r)$ for every law; put
$Y=\min(1,e(1-U)/s)$ for $s>0$. The maximal absent weight is $lY$,
and $a=l$ is still qualified. At either zero-probability
history the multiplied inequalities impose only the constraint
appropriate to that history. Substitution yields the linear
objective
\[
        (1-l)(1-s)x+l\big[L+(1-s-L)Y-k\big].
\]
The optimal choice is an endpoint, with the no-probe endpoint
canonical at equality. The difference in optimal fraction is a
difference of risk doctrine, not an inferred change in RF.

For $s=d=1/5$, $e=1/10$, $k=1/20$, the exact whole-plan optimum
uses the probe in half of trials and obtains $27/40$.
The observation-conditional optimum probes every trial, obtaining
$57/100$. For the larger next duty, no probe is optimal in both
doctrines and obtains $2/5$. To see dominance for that size,
never admit on the costly positive outcome (it always fails).
Reassign the weighted absent admissions and no-probe admissions
to the no-probe branch. Its chance of a good callback cannot
decrease for any law; global risk remains $s(n+b)$.
Under the conditional rules, $n+b\le x$ so that same no-probe
policy is also qualified. Information may be real and still
destroy the optional service it was bought to improve.

\paragraph{Comparators and limits.}
The exact rational direct polytope uses the four atoms and their
actual outcome tuples; endpoints suffice for its affine robust
constraints. It matches the analytical policy objective and least
optimal probe fraction on 400 benchmarks plus the named example.
Three Lean kernels verify the objective envelope and fraction bound.
The work distinguishes attained receipt-conditioned capacity from
independent law uncertainty. It does not justify physical drop
marginals, a change in authorized risk semantics, a different
available action set, deployed control, or independent mission
selection. An equal-information direct planner remains a tie.
\end{document}
